\section{研究前沿与开放问题}

\subsection{理论与实践之间仍有明显断层}
讲者明确指出，当前主流后训练算法与理想化机器学习理论并未完全对齐。实践中很多成功经验来自规模化试验和工程迭代，而不是闭式推导。因此，未来几年算法空间仍非常开放。

\begin{importantbox}{开放研究方向}
\begin{itemize}
    \item 更贴合 Agent 场景的策略优化目标；
    \item 更稳健的长轨迹 credit assignment；
    \item 更低成本的高质量 verifier 构建；
    \item 训练与部署一体化的在线反馈闭环。
\end{itemize}
\end{importantbox}

\subsection{从人类学习类比到可实现算法}
课程多次使用 ``人类学习'' 做类比：先看示例，再做题，再拿反馈迭代。类比本身直观，但真正困难在于把它转成可训练、可扩展、可审计的算法与系统。这一点在 Agent 场景尤为突出。

\begin{knowledgebox}{类比可用，直译不可用}
把人类学习过程直接翻译成算法通常会失败。可行路径是提取可操作原则，例如：
\begin{itemize}
    \item 先建立基本可行策略（Light SFT）；
    \item 再以高质量反馈持续探索（RL）；
    \item 始终监控泛化与评估污染（Holistic eval）。
\end{itemize}
\end{knowledgebox}

\subsection{部署风险与治理议题}
当 Agent 被用于代码、业务流程或关键系统，训练问题会直接转化为治理问题。模型是否会越权调用、是否会在弱验证场景投机、是否能被追责审计，都需要在训练和评估阶段前置设计。

\begin{warningbox}{能力提升必须伴随可控性提升}
如果只提升任务成功率而不提升可解释性、可回滚性和权限治理，系统风险会随能力线性以上增长。可验证训练的真正价值，正在于把能力增长与治理能力绑定在一起。
\end{warningbox}

\subsection{本章小结}
本章结论是：可验证 Agent 仍处于快速演化期，方法论尚未收敛。最值得投入的方向是把 ``高质量反馈学习'' 与 ``可部署治理机制'' 联合起来。

